\subsection{Digital Prototype}\label{subsect:DigitalPrototype}
% Link o QR code al progetto Figma interattivo (sul modello del QR code in
% Hammami_Hub, fig. 3.1). Intro breve: da qui in poi i mockup digitali sono
% raggruppati per flusso (non una schermata per sottosezione), sul modello di
% BookIT 4.4 / Hammami_Hub 3.4 -- ogni flusso ha gli screenshot delle
% schermate coinvolte in sequenza, e un diagramma mermaid a frecce quando il
% flusso dirama (piu' di un percorso possibile) o le transizioni non sono
% ovvie dal solo screenshot. Tre gruppi: comuni ai due ruoli, poi cliente,
% poi professionista.

\paragraph{Common Flows}

\subsubsection{Login \& Reset Password Flow}\label{subsubsect:LoginResetFlow}
% Login -> Forgot Password -> Reset Password, piu' il redirect post-login al
% ruolo corretto (/m o /p, in base al profilo). Wireframe esistenti solo in
% doc/assets/pt/01-03 (Login, Forgot Password, Set Password) -- condivisi tra
% i due ruoli, nessun set separato lato cliente. Diagramma mermaid lineare,
% con un ramo per "password dimenticata" e uno per il login diretto.

\subsubsection{Plan Creation \& Editing Flow}\label{subsubsect:PlanCreationEditingFlow}
% CreatePlanFlow.jsx e' condiviso testualmente tra i due ruoli ("Used by both
% roles. Only memberId, authorId and onDone differ.") -- un solo flusso, non
% duplicato. Ramo principale: 2 step in-pagina (nome piano -> prima sessione
% con esercizi aggiunti via autocomplete inline), nessuna schermata di
% riepilogo -- torna dritto alla panoramica del piano. Entry point diverso
% per ruolo: tab Workout per il cliente, Client Detail per il professionista
% -- da annotare nel diagramma, non da raccontare due volte.
%
% Ramo secondario, editing di un piano gia' esistente (aggiungere una
% sessione o un esercizio): stesso SessionForm, ma interfaccia diversa per
% ruolo -- lato cliente sono due schermate a parte (Add a Session, Add an
% Exercise), lato professionista e' tutto inline dentro Client Detail /
% Workout Plan (nessuna navigazione). Il diagramma deve mostrare questa
% divergenza, non solo il ramo comune.

\paragraph{Client Flows}

\subsubsection{Home \& Global Navigation Flow}\label{subsubsect:ClientHomeNavFlow}
% Home come screenshot d'apertura della sezione Cliente. Frecce etichettate
% "tab bar" verso Workout, Nutrition, My Trainer; freccia etichettata
% "avatar" verso Profile. Sono i punti d'accesso sempre raggiunti dalla shell
% dell'app (bottom nav + top bar), non link dentro il contenuto della Home --
% da leggere come premessa visiva ai flussi seguenti, non come schermata con
% link propri.

\subsubsection{Workout Plan Exploration Flow}\label{subsubsect:WorkoutExplorationFlow}
% Workout Plan -> Session Detail -> Exercise Detail, senza avviare una
% sessione live. Percorso di sola consultazione.

\subsubsection{Live Session Flow}\label{subsubsect:LiveSessionFlow}
% Session Detail -> Live Session -> Exercise Detail Live (stato inline dentro
% Live Session, non una route separata) -> Open-Run Sheet / End-Run Sheet
% (overlay) -> Congratulations Dialog (overlay) -> Session Summary. Diagramma
% mermaid con frecce per le diramazioni: Open-Run Sheet quando un secondo run
% si apre mentre uno e' gia' attivo, "Stopped Early" quando si esce da
% End-Run Sheet senza completare.

\subsubsection{Nutrition Flow}\label{subsubsect:NutritionFlow}
% Nutrition Plan -> Meal Detail (route /m/nutrition/meal/:mealId gia'
% esistente, wireframe ancora da fare -- nodo comunque presente, non
% taciuto). Ramo stato vuoto (nessun piano assegnato) -> Book Appointment,
% che rimanda al Personal Trainer \& Booking Flow.

\subsubsection{Personal Trainer \& Booking Flow}\label{subsubsect:TrainerBookingFlow}
% My Trainer -> Chat; My Trainer -> View Appointments -> Booking Sheet ->
% conferma prenotazione. Ramo "nessun professionista assegnato" -> Browse
% Trainers (redirect automatico da My Trainer quando non c'e' ancora un
% professionista scelto).

\subsubsection{Profile Flow}\label{subsubsect:ClientProfileFlow}
% Profile -> Access Badge, Subscription, Rewards, Settings. Diagramma a
% stella: tutti i rami partono da Profile, nessun collegamento tra loro.

\paragraph{Professional Flows}

\subsubsection{Home \& Global Navigation Flow}\label{subsubsect:ProHomeNavFlow}
% Stessa struttura del flusso Home lato cliente: Home come screenshot
% d'apertura, frecce "tab bar" verso Clients, Calendar, Chat; freccia
% "avatar" verso Profile. L'agenda mostrata in Home include anche
% un'anteprima appuntamenti che rimanda al Calendar Flow -- non ripetuta qui,
% solo richiamata.

\subsubsection{Calendar Flow}\label{subsubsect:CalendarFlow}
% Calendar -> Availability. Calendar -> Appointment Detail (raggiungibile
% anche dall'anteprima appuntamenti in Home). Nuovo appuntamento ("+") come
% azione inline sulla stessa schermata di Calendar, non una route a parte.

\subsubsection{Client Management Flow}\label{subsubsect:ClientManagementFlow}
% Clients -> Client Detail -> quattro rami: Chat (stesso nodo Chat della tab
% bar -- doppia via d'accesso, un solo screenshot di destinazione); Workout
% Plan (rimanda al Plan Creation \& Editing Flow comune, meccanismo identico
% a quello del cliente); Nutrition Plan (consultazione + creazione/editing
% inline -- ramo professionista-only, il cliente non ha mai un bottone "+"
% qui); Progress Tracking (sola consultazione, nessuna creazione).

\subsubsection{Profile \& Settings Flow}\label{subsubsect:ProProfileSettingsFlow}
% Profile -> Settings. Molto piu' corto della versione cliente: nessun badge,
% subscription o rewards lato professionista.
